\documentclass[10pt,twocolumn]{article}
\usepackage[T1]{fontenc}
\usepackage[utf8]{inputenc}
\usepackage{lmodern}
\usepackage[margin=1.8cm,top=2.4cm,bottom=2cm,columnsep=0.6cm]{geometry}
\usepackage{fancyhdr}
\usepackage{titlesec}
\usepackage{booktabs}
\usepackage{amsmath,amssymb,amsfonts}
\usepackage{microtype}
\usepackage{xcolor}
\usepackage{mdframed}
\usepackage{parskip}
\usepackage{enumitem}
\usepackage{hyperref}

\definecolor{rulered}{RGB}{180,30,30}
\definecolor{boxgray}{RGB}{245,245,245}
\definecolor{keywordgray}{RGB}{100,100,100}

\pagestyle{fancy}
\fancyhf{}
\fancyhead[L]{\textsc{\small Claude Mag}}
\fancyhead[C]{\small\textit{Machine Learning Research Digest}}
\fancyhead[R]{\small 2026-08-05}
\fancyfoot[C]{\small\thepage}
\renewcommand{\headrulewidth}{0.8pt}

\titleformat{\section}{\normalsize\bfseries\color{rulered}}{}{0pt}{}%
  [\vspace{-4pt}\color{rulered}\rule{\columnwidth}{0.4pt}]
\titlespacing{\section}{0pt}{8pt}{4pt}

\hypersetup{colorlinks=true,urlcolor=rulered,linkcolor=black}

\begin{document}
\setlength{\parindent}{0pt}
\setlength{\parskip}{4pt}

{\small\color{keywordgray}\textbf{Keywords:}
  diffusion transformer \textbullet{}
  momentum orthogonalization \textbullet{}
  optimizer \textbullet{}
  image generation}\par\vspace{4pt}

{\LARGE\bfseries\raggedright
  CMuon: Accelerating and Stabilizing Diffusion Transformer Training
  via Chunked Momentum Orthogonalization}\par\vspace{4pt}

{\large\itshape\raggedright
  Partitioning Fused Weight Matrices Before Orthogonalization Delivers a $2\times$
  Speedup Over AdamW and FID 1.18 on ImageNet~256}\par\vspace{6pt}

{\small\textbf{By} Chuyan Chen, Peng Sun \& Kun Yuan}\par\vspace{2pt}

{\small\color{keywordgray}arXiv:2608.02502 \textbullet{} ECCV 2026}
\par\vspace{2pt}\noindent{\color{rulered}\rule{\columnwidth}{0.6pt}}\vspace{6pt}

The Muon optimizer's momentum orthogonalization accelerates standard Transformer
training, but applying it directly to Diffusion Transformers stalls at late stages.
CMuon traces the cause to fused weight tensors coupling independent subspaces and
fixes it with a single structural change---partitioning matrices into chunks before
orthogonalization. A 675M-parameter DiT trained with CMuon achieves FID~1.18 on
ImageNet 256$\times$256 in 200 epochs, more than $2\times$ faster than AdamW.

\begin{mdframed}[backgroundcolor=boxgray,linecolor=rulered,linewidth=1pt,
    innerleftmargin=6pt,innerrightmargin=6pt,
    innertopmargin=6pt,innerbottommargin=6pt]
\textbf{\small AT A GLANCE}\par\vspace{4pt}
\begin{description}[leftmargin=0pt,labelsep=4pt,itemsep=2pt,parsep=0pt,
    font=\small\bfseries]
  \item[Problem:] Muon achieves faster convergence than AdamW on standard Transformers
    but stalls in the late training phase when applied to Diffusion Transformers (DiTs).
  \item[Why it matters:] DiT training is the dominant compute cost for frontier image
    and video generation; a $2\times$ optimizer speedup directly halves training cost.
  \item[Method:] Identify fused weight tensors (AdaLN, QKV) as the source of subspace
    coupling; chunk them prior to Nesterov momentum orthogonalization.
  \item[Result:] FID 1.18 on ImageNet 256$\times$256 in 200 epochs (675M parameters),
    $>2\times$ training speedup over AdamW.
  \item[Evidence:] Controlled ablations against Muon and AdamW; FID curves over full
    200-epoch training run; accepted ECCV 2026.
\end{description}
\end{mdframed}

\section{The Big Picture}

Diffusion Transformers (DiTs) power the best image and video generation models of
2026, from FLUX to Sora-class video synthesis. Training them is extremely expensive:
a state-of-the-art image generation DiT requires thousands of GPU-hours, and video
models demand an order of magnitude more. Any optimizer improvement that accelerates
convergence without degrading final quality translates directly to lower training cost.

Muon---the Momentum Orthogonalization optimizer---has attracted attention because it
achieves significantly faster convergence than AdamW on vanilla Transformer language
models. CMuon makes Muon work on DiTs by identifying and fixing a structural mismatch
between Muon's assumptions and DiT's weight layout.

\section{Why Existing Approaches Fall Short}

\textbf{AdamW} is the de-facto DiT optimizer. It works reliably but converges slowly,
requiring many epochs to reach competitive FID scores.

\textbf{Vanilla Muon} applies Nesterov momentum and then orthogonalizes the resulting
gradient update matrix via Newton--Schulz iterations before applying the weight update.
This exploits the geometry of the weight matrix space to find descent directions that
make maximal use of each parameter. On standard Transformer blocks (Q, K, V, and MLP
projections each in their own tensors), it works well.

But DiTs use \textbf{Adaptive Layer Norm (AdaLN)}: a conditioning mechanism that
generates per-channel scale and shift parameters from the diffusion timestep embedding,
modulating the activations of each Transformer block. For GPU efficiency, the linear
layers implementing AdaLN fuse the weights for scale and shift into a single large
matrix, and similarly the Q, K, V projections are sometimes stored fused. Applying
Muon to these fused tensors treats two functionally independent linear maps as a single
entity, causing the orthogonalization to mix their update directions. This implicit
coupling distorts the gradient geometry and causes convergence to plateau.

\section{Core Mechanism}

CMuon introduces a minimal pre-processing step before the Muon update:

\begin{enumerate}[itemsep=2pt,parsep=0pt]
  \item \textbf{Identify} all fused weight tensors in the DiT (AdaLN scale+shift
    matrices, fused QKV projections, etc.).
  \item \textbf{Chunk} each fused matrix along its output channel dimension according
    to the functional boundary (e.g., split AdaLN into its scale matrix and its shift
    matrix separately).
  \item \textbf{Orthogonalize} each chunk independently via the Newton--Schulz
    approximation to the matrix orthogonal factor.
  \item \textbf{Reassemble} the chunked updates and apply them to the original
    fused parameter tensor.
\end{enumerate}

No architectural changes or additional hyperparameters are required. The same chunking
boundaries used during the forward pass for functional interpretation are reused here
during the optimizer step.

\section{Technical Formulation}

Let $\mathbf{W} \in \mathbb{R}^{2d \times d}$ be a fused AdaLN weight matrix
combining a scale projection $\mathbf{W}_s \in \mathbb{R}^{d \times d}$ and a
shift projection $\mathbf{W}_b \in \mathbb{R}^{d \times d}$.

The Muon update for $\mathbf{W}$ is:
\[
\Delta\mathbf{W} = \mathrm{NS}(\mathbf{M}),
\quad
\mathbf{M} = \mu \mathbf{M}_{\mathrm{prev}} + \mathbf{G},
\]
where $\mathbf{G}$ is the gradient, $\mu$ is momentum, and
$\mathrm{NS}(\cdot)$ denotes the Newton--Schulz orthogonalization:
\[
\mathrm{NS}(\mathbf{A}) \approx \frac{3}{2}\mathbf{A} - \frac{1}{2}\mathbf{A}\mathbf{A}^\top\mathbf{A}
\quad \text{(one step)}.
\]

The CMuon update orthogonalizes each chunk separately:
\[
\Delta\mathbf{W}_s = \mathrm{NS}(\mathbf{M}_s), \quad
\Delta\mathbf{W}_b = \mathrm{NS}(\mathbf{M}_b),
\]
where $\mathbf{M}_s, \mathbf{M}_b$ are the momentum matrices for each chunk
(maintained separately). The final update is $\Delta\mathbf{W} =
[\Delta\mathbf{W}_s;\,\Delta\mathbf{W}_b]$.

By orthogonalizing each functional subspace independently, each chunk's update
direction is computed in the geometry appropriate to its own weight space rather
than the mixed geometry of the fused tensor.

\section{Learning or Inference Procedure}

CMuon is a drop-in optimizer replacement. Training proceeds as:
\begin{enumerate}[itemsep=2pt,parsep=0pt]
  \item Forward pass and diffusion denoising loss computation (identical to standard
    DiT training).
  \item Backward pass yielding gradients for all parameters.
  \item For each fused weight tensor, split gradients and momentum along the chunk
    boundary, orthogonalize each chunk, reassemble update.
  \item For non-fused parameters (e.g., biases, embedding tables), use AdamW as usual
    (Muon is not well-suited to 1-D parameter vectors).
  \item Apply updates with a learning rate schedule identical to AdamW baselines
    for fair comparison.
\end{enumerate}

Inference is unaffected: CMuon is a training-time optimizer; the deployed model
is identical in architecture to a standard DiT.

\section{What the Guarantee Says}

Muon's orthogonalization is motivated by the Riemannian gradient on the manifold of
matrices with bounded spectral norm. Each update step is guaranteed to be a descent
direction in this geometry. The CMuon insight is that applying this guarantee to a
fused matrix provides a weaker guarantee than applying it to each component separately,
because the fused matrix does not lie on a single manifold coherent with either
component's optimization landscape. Chunk-level orthogonalization restores the
per-component descent guarantee.

\section{Experimental Findings}

Evaluated on DiT-XL/2 (675M parameters) training on ImageNet 256$\times$256
with class-conditional generation:

\begin{itemize}[itemsep=2pt,parsep=0pt]
  \item CMuon achieves \textbf{FID 1.18} in 200 epochs---competitive with the
    best published DiT-XL/2 results.
  \item Training is \textbf{$>2\times$ faster} than AdamW measured by FID at equal
    epoch count.
  \item Vanilla Muon on DiTs plateaus around epoch 120--140 and never reaches FID~1.3;
    CMuon continues descending through 200 epochs.
  \item No guidance (classifier-free or otherwise) is used during evaluation,
    making this an apples-to-apples comparison.
\end{itemize}

\section{Ablations and Interpretation}

\textbf{Without chunking (vanilla Muon):} Convergence stalls at FID $\approx 1.4$
around epoch 130. The late-stage plateau is eliminated by chunking.

\textbf{Chunk granularity:} Chunking at the AdaLN scale+shift boundary is sufficient.
Further splitting within Q, K, V provides marginal additional gain on this architecture.

\textbf{Non-fused layers:} Applying Muon (without chunking) to non-fused DiT layers
(attention projections stored separately) does not cause convergence issues, confirming
that fused tensors are the specific failure mode.

\textbf{AdamW vs.\ CMuon FID curves:} At epoch~100 AdamW achieves FID~$\approx$1.8,
while CMuon achieves FID~$\approx$1.3, confirming the $2\times$ epoch-efficiency claim.

\section*{Reference}

Chuyan Chen, Peng Sun, Kun Yuan. \textbf{CMuon: Accelerating and Stabilizing
Diffusion Transformer Training via Chunked Momentum Orthogonalization}.
arXiv:2608.02502, ECCV 2026.
\url{https://arxiv.org/abs/2608.02502}

\end{document}
